%%%%%%%%%%%%Outline%%%%%%%%%%%%%%%%%%

% section 3 Motivation %
% message<eBPF optimization is constrained by two fundamental design problems: rich optimization logic should not live on the kernel’s safety-critical compilation path, and optimization profitability cannot be decided once and for all at first load.> %
% info 1 Closing the current code-quality gap requires more than one missing rewrite; it requires a growing optimizer. %
% info 2 Putting that optimizer into the kernel would impose both high engineering cost and higher-consequence security risk. %
% info 3 Even if richer optimization logic existed, a fixed load-time policy would still be wrong for some deployments because profitability depends on runtime behavior. %

% section 3 opening paragraph %
% message<This section motivates a split of responsibilities: the kernel should keep only a small trusted recompilation mechanism, while userspace should host rich optimization logic and make optimization decisions dynamically after deployment.> %
% info 1 The paragraph should stand alone as the thesis of Section 3, without referring back to Section 2. %
% info 2 It should preview the two subsection claims in order: placement first, timing second. %
% info 3 It should end by foreshadowing the design consequence: kernel safety enforcement plus userspace policy. %

% section 3.1 Rich Optimization Logic Belongs Outside the Kernel %
% paragraph 1 %
% message<The kernel already has a natural safety boundary for BPF execution: the verifier.> %
% info 1 The verifier checks memory safety, bounded execution, control-flow validity, and type-correct interaction with kernel objects. %
% info 2 Only programs that pass verification may execute in kernel context. %
% info 3 A rewritten program can therefore be treated as a fresh candidate program and re-verified before installation. %

% paragraph 2 %
% message<There is real code-quality headroom in the current kernel JIT, but the gap is broad rather than isolated.> %
% info 1 A hash-style example should illustrate missed rotate, extract, endian, or byte-recomposition opportunities. %
% info 2 A binary-search-style example should illustrate missed conditional-move or branch-layout opportunities. %
% info 3 The purpose of these two examples is to show that the problem is not one corner case, but multiple recurring classes of missed optimization. %

% paragraph 3 %
% message<Closing this gap requires a growing collection of passes rather than a few isolated backend tweaks.> %
% info 1 One rewrite can repair one pattern, but sustained code-quality improvement needs multiple analyses, transforms, and architecture-specific lowerings. %
% info 2 Different optimization families solve different recurring misses and interact with one another. %
% info 3 Once the goal is broad code-quality improvement, the implementation starts to resemble an optimizer ecosystem rather than a small JIT backend. %

% table 1 %
% message<Engineering-effort comparison between kernel-side BPF transformation logic, userspace optimization passes, and representative LLVM pass implementations.> %
% info 1 Rows should include representative LLVM passes, kernel BPF verifier/JIT transformation code, and BpfReJIT userspace pass code. %
% info 2 Columns should include component, privilege boundary, approximate LOC/SLOC, and role. %
% info 3 Purpose: show that rich optimization logic is multi-thousand-line engineering and therefore a poor fit for the trusted in-kernel compilation path. %

% paragraph 4 %
% message<The in-kernel approach is unattractive because the engineering effort is large, persistent, and accumulative.> %
% info 1 Rich optimization logic is not a one-time patch; it requires iteration, retuning, testing, and maintenance. %
% info 2 Kernel-resident logic must be reviewed across architectures and carried across kernel releases. %
% info 3 This explains why code quality remains limited: the kernel is a poor home for a steadily expanding pass ecosystem. %

% paragraph 5 %
% message<Putting more optimization logic in the kernel also changes the consequence of mistakes.> %
% info 1 A userspace optimizer bug is still serious, but it remains an ordinary userspace fault. %
% info 2 A verifier- or JIT-path bug sits on the privilege boundary for untrusted BPF execution. %
% info 3 As more optimization logic moves into that path, the system accumulates more security-sensitive kernel complexity. %

% figure 1 %
% message<Security-survey figure supporting the claim that expanding privileged compiler logic in the kernel is risky.> %
% info 1 Preferred form: a timeline or histogram of public BPF verifier/JIT-related vulnerabilities, grouped by year and component. %
% info 2 Optional secondary series: a userspace-compiler bug-count reference only if the methodology is truly comparable and defensible. %
% info 3 Purpose: support the three core claims of 3.1: engineering effort is high, extra kernel logic increases security-sensitive complexity, and failures in that path are more severe in consequence. %

% paragraph 6 %
% message<The right design consequence is to keep only a small recompilation mechanism in the kernel and place rich optimization logic in userspace.> %
% info 1 The kernel should recover program form, re-verify rewrites, and safely replace code. %
% info 2 Userspace should host the pass pipeline, architecture-aware rewrite logic, and rapid iteration. %
% info 3 This preserves the verifier safety boundary without giving up optimization capability. %

% section 3.2 Optimization Profitability Is a Post-Deployment Property %
% paragraph 1 %
% message<Even with richer optimization logic, a fixed static policy is still the wrong abstraction because profitability is not a load-time property.> %
% info 1 Legality can be decided at load time. %
% info 2 Profitability depends on runtime behavior that is not yet visible at load time. %
% info 3 A fixed load-time optimizer therefore freezes the decision too early. %

% paragraph 2 %
% message<The same legal optimization can speed up some programs and slow down others.> %
% info 1 One group of benchmarks should show clear wins from a single pass. %
% info 2 Another group should show clear regressions from that same pass. %
% info 3 The message is that legality does not imply profitability. %

% paragraph 3 %
% message<Fixed combinations of individually reasonable passes still do not dominate across programs.> %
% info 1 Interactions through code layout, register pressure, and hot-path reshaping can create regressions. %
% info 2 The problem is therefore deeper than choosing one better fixed heuristic. %
% info 3 This is why the current Table 3 and Table 4 story should be compressed into one visual argument instead of staying as separate tables. %

% figure 2 %
% message<Two-panel policy-sensitivity figure showing why static optimization policy fails.> %
% info 1 Panel (a): one pass across several benchmarks, with both positive and negative bars. %
% info 2 Panel (b): a fixed-all or multi-pass configuration across several benchmarks, again with both wins and regressions. %
% info 3 Purpose: show visually that no single static optimization policy dominates across deployed workloads. %

% paragraph 4 %
% message<These mixed outcomes are expected because runtime behavior, not just bytecode shape, determines profitability.> %
% info 1 Branch predictability depends on live input distributions. %
% info 2 Hot paths depend on deployment behavior after program load. %
% info 3 Workload skew, value stability, and machine-specific effects are invisible to purely static load-time analysis. %

% paragraph 5 %
% message<Optimization policy must therefore remain dynamic after deployment, while the kernel continues to decide safety.> %
% info 1 Already-loaded programs must remain revisit-able after load. %
% info 2 Rewrite selection must be able to use runtime signals from the actual deployment. %
% info 3 This leads directly to a post-load system with userspace policy and kernel safety enforcement. %

%%%%%%%%%%%%Outline%%%%%%%%%%%%%%%%%%

\section{Motivation}
\label{sec:motivation}

eBPF needs better native code, but neither a richer in-kernel optimizer nor a fixed first-load policy is the right way to get it. Closing the current code-quality gap requires a growing collection of analyses, rewrites, and architecture-specific lowerings rather than one missing JIT tweak. Putting that optimizer on the verifier and JIT path would substantially increase both engineering cost and the consequence of mistakes in one of the kernel's most security-sensitive code paths. Even if that engineering cost were acceptable, a fixed load-time policy would still make the wrong choice for some programs, because optimization profitability depends on runtime behavior that appears only after deployment. This section therefore motivates a different split of responsibilities: the kernel should keep only a small trusted recompilation mechanism, while userspace should host rich optimization logic and make optimization decisions dynamically after deployment.

\subsection{Keeping the Kernel Compilation Path Small}
\label{sec:motivation-kernel}

The kernel already has a natural safety boundary for BPF execution: the verifier~\cite{verifier}. Before a program may execute in kernel context, the verifier checks memory safety, bounded execution, control-flow validity, and type-correct interaction with kernel objects. That boundary already answers the safety question for optimization. A rewritten program can be treated as another candidate BPF program and re-verified before installation, so the kernel does not need to trust an external optimizer's policy. It only needs to keep the authority to accept or reject candidate rewrites.

There is real code-quality headroom in the current kernel JIT, but the gap is broad rather than isolated. A hash-style kernel exposes missed opportunities such as rotate, extract, endian, and byte-recomposition simplification. A branch-heavy binary search exposes missed opportunities such as conditional-move selection and branch-layout improvement. These examples matter because they are ordinary optimization families rather than one-off corner cases. The code-quality problem is therefore not that the kernel misses one clever rewrite. It is that multiple recurring classes of missed optimization remain visible in common code patterns.

Closing this gap requires a growing collection of passes rather than a few isolated backend tweaks. A rotate or extract pass does not solve byte recomposition. A byte-recomposition pass does not solve conditional-move lowering. A conditional-move pass still leaves cleanup, profitability, and architecture-specific lowering decisions unresolved. Once the goal is sustained code-quality improvement, the implementation starts to resemble an optimizer ecosystem rather than a small JIT backend. Table~\ref{tab:motivation-effort} is therefore part of the argument rather than a side detail: it quantifies the engineering burden of representative LLVM passes, kernel-side BPF transformation logic, and \tool's userspace pass code.

\begin{table}[t]
  \centering
  \caption{Placeholder for the engineering-effort comparison in \S\ref{sec:motivation-kernel}. The final table will compare representative LLVM pass implementations, current kernel-side BPF transformation/JIT code, and \tool's userspace pass code. Its purpose is to show that rich optimization logic is already multi-thousand-line engineering and is therefore a poor fit for the trusted in-kernel compilation path.}
  \label{tab:motivation-effort}
  \small
  \fbox{\parbox{0.92\linewidth}{\centering
    Placeholder for Table~\ref{tab:motivation-effort}: \\
    component \quad privilege boundary \quad approximate LOC/SLOC \quad role}}
\end{table}

The in-kernel approach is unattractive because the engineering effort is large, persistent, and accumulative. Rich optimization logic is not a one-time patch. It must be extended, retuned, tested, and maintained across processor families and kernel releases. Kernel-resident logic also carries a long review tail, because each added optimization must remain compatible with verifier behavior, JIT backends, and architecture-specific corner cases. This is exactly why code quality remains limited in the first place: the kernel is a poor home for a steadily expanding pass ecosystem.

Putting more optimization logic in the kernel also changes the consequence of mistakes. A userspace optimizer bug is still serious, but it remains an ordinary userspace fault. A verifier- or JIT-path bug sits on the privilege boundary for untrusted BPF execution, where the cost of being wrong is much higher. As more optimization logic moves into that path, the system accumulates more security-sensitive kernel complexity rather than merely more compiler code. Figure~\ref{fig:motivation-security} will summarize this point with a survey of public BPF verifier/JIT security issues.

\begin{figure}[t]
  \centering
  \fbox{\parbox{0.92\linewidth}{\centering
    Placeholder for Figure~\ref{fig:motivation-security}: \\
    security-survey figure for BPF verifier/JIT issues by year and component}}
  \caption{Placeholder for a security-survey figure. The final figure will summarize public BPF verifier/JIT-related security issues, ideally grouped by year and component. If a userspace-compiler reference series turns out to be methodologically comparable, it can be added as a secondary series. The purpose of the figure is to support the three core claims of \S\ref{sec:motivation-kernel}: rich in-kernel optimization logic is expensive to engineer, increases security-sensitive complexity, and makes mistakes more severe in consequence.}
  \label{fig:motivation-security}
\end{figure}

The right design consequence is therefore to keep only a small recompilation mechanism in the kernel and place rich optimization logic in userspace. The kernel should recover program form, re-verify rewrites, and safely replace code if verification succeeds. Userspace should host the pass pipeline, architecture-aware rewrite logic, and rapid iteration. This split preserves the verifier's safety boundary without giving up the richer optimization logic that better code quality requires.

\subsection{Optimization Profitability Is a Post-Deployment Property}
\label{sec:motivation-dynamic}

Even with richer optimization logic, a fixed static policy is still the wrong abstraction because profitability is not a load-time property. Legality can be decided at load time, but profitability often cannot. The properties that determine whether a rewrite helps---branch predictability, hot-path frequency, value stability, traffic mix, and machine-specific behavior---appear only after deployment. A fixed load-time optimizer therefore freezes the decision too early, before the most relevant information exists.

The same legal optimization can speed up some programs and slow down others. A conditional-move rewrite is a simple example. It helps when branch outcomes are hard to predict, but it can lose when those same branches are already highly predictable~\cite{cmov_branch}. The point is not a special case for one instruction. The point is that legality does not imply profitability: a pass may be correct for every program and still be beneficial for only some of them.

Fixed combinations of individually reasonable passes do not solve this problem. Each rewrite may look attractive in isolation, yet their interaction can still create regressions through code layout, register pressure, and hot-path reshaping. The issue is therefore deeper than finding one better fixed heuristic. What fails is the assumption that there exists a single static optimization policy that reliably dominates across deployed programs. Figure~\ref{fig:motivation-policy} will compress this argument into one visual comparison: one panel for a single optimization that produces both wins and regressions, and one panel for a fixed multi-pass configuration that still produces both wins and regressions.

\begin{figure*}[t]
  \centering
  \fbox{\parbox{0.95\linewidth}{\centering
    Placeholder for Figure~\ref{fig:motivation-policy}: two-panel bar chart. \\
    Panel (a): one pass across several benchmarks, with both wins and regressions. \\
    Panel (b): fixed-all or multi-pass configuration across several benchmarks, again with both wins and regressions.}}
  \caption{Placeholder for the policy-sensitivity figure. The final figure will replace the current separate tables with one two-panel visual argument. Panel~(a) will show that the same legal pass improves some programs and hurts others. Panel~(b) will show that a fixed combination of individually reasonable passes still fails to dominate across programs. Together, the two panels motivate dynamic post-load optimization rather than a fixed static policy.}
  \label{fig:motivation-policy}
\end{figure*}

These mixed outcomes are expected because runtime behavior, not just bytecode shape, determines profitability. Branch predictability depends on live input distributions. Hot paths depend on deployment behavior after program load. Workload skew, value stability, and machine-specific effects are invisible to purely static load-time analysis. Load-time optimization can inspect code and loader context, but it cannot observe the steady-state behavior that often determines whether a rewrite is actually worthwhile.

Optimization policy must therefore remain dynamic after deployment, while the kernel continues to decide safety. Already-loaded programs must remain revisit-able after load, and rewrite selection must be able to use runtime signals from the actual deployment. This requirement is separate from safety rather than in conflict with it: the kernel still decides whether a rewritten program is legal, while userspace decides whether that legal rewrite is worth trying now. The next section presents \tool as a post-load system built around exactly this division of labor.